\section{The Quantum Sector}
\label{sec:quantum}

The quantum pillars are where the model is most careful, because a smooth complex amplitude cannot
live on a three-valued vibe. This section walks the quantum findings (P6, P31, P70).

\paragraph{Amplitudes as patch averages.}
A quantum amplitude is not carried by a single vibe. It is an average over a patch of many vibes:
its magnitude is a density of how many are excited together, its phase the patch's shared rhythm.
This is the dictionary that lets a continuous wavefunction emerge from the discrete substrate, and
it is what makes the next results possible.

\paragraph{The pillars (P6, P31).}
With that dictionary, the quantum pillars appear. Evolution conserves probability, and a disturbance
spreads in the wave-like way that signals coherence. Interference is present, possibilities adding
and cancelling. The Bell correlations are reproduced, and the model says why: because there is no
hidden state, the measuring device and the measured system are made of the same mesh and share a
common past, so they are not independent, and the shared substrate carries the correlation with no
signal travelling faster than light. P31 assembles the quantum formalism on the mesh, and P6
establishes the Bell result.

\paragraph{The Born rule, derived (P70).}
P70 derives the Born rule rather than postulating it. The amplitude's magnitude is the square root
of a density of vibes, so fair sampling of the substrate gives the squared amplitude, and the
exponent two is singled out. The testbed shows this two ways, a counting argument and an
envariance argument \citep{zurek2005probabilities}, both returning $|c|^2$, with the powers $p = 1$ and $p = 3$
excluded.

 So probability equal to amplitude squared is a consequence of the substrate's structure,
not an axiom.

\paragraph{Contact with data (P35).}
P35 makes the order-of-magnitude contact with experiment that P82 later sharpens: the model's
distinctive effects sit within the observed bounds, and the everpresent cosmological constant of
Section~\ref{sec:cosmology} matches the observed dark-energy scale to order of magnitude.
